\documentclass[11pt,a4paper]{article}
\usepackage[margin=23mm]{geometry}
\usepackage{amsmath,amssymb,amsthm}
\newtheorem{theorem}{Theorem}
\newcommand{\gen}[1]{\langle#1\rangle}
\title{The tail after a critical head is intransitive}
\author{Human Compression project: L044 continuation}
\date{}
\begin{document}\maketitle
Retain all acquired critical-head hypotheses and proofs, including
the smaller-degree alternating induction used to reach that
configuration. Thus $G\le A_n$ has a transitive connected string
C-marking, all full maximal omissions are intransitive, and rank
failure has already given $n=2r\ge18$. The first three marks
form the actual six-point head on $D$, and with
$E=\Omega\setminus D$ one has
\[
 \rho_3=(6\ b)\delta,\quad
 \operatorname{supp}(\delta)\subseteq E\setminus\{b\},\qquad
 B=G_{\ge4},\quad B|_D=1.
\]
All generators of $B$ except its first fix $b$.

\begin{theorem}
The group $B$ is intransitive on $E$. Consequently colour $3$
has at least two full ambient-eligible edges.
\end{theorem}
\begin{proof}
Suppose $B$ is transitive. Restriction to $E$ is faithful,
since every element fixes $D$. It therefore transfers the
full intersection property, independence and the unchanged
nonidentity marks. Its degree and rank are
\[
                 f=|E|=n-6\ge12,\qquad
                 d=r-4=f/2-1.
\]
Every mark is even on $E$. Its noncommutation diagram is
connected, as a terminal interval of the connected full string.
Every distinguished omission in $B$ is intransitive on $E$:
otherwise the head and the edge $(6\ b)$ of $\rho_3$
would make the corresponding full ambient omission transitive.

If $B$ is proper primitive in $A_E$, the complete accepted
primitive bound gives
\[
 d\le\left\lfloor\frac{f-3}{2}\right\rfloor=f/2-2,
\]
contrary to its displayed rank. This uses the full primitive
proof and its required finite boundary; it is not an invocation
of the alternating rank theorem at degree $f$.
If $B=A_E$ naturally, the accepted pure-head obstruction
applies directly, with $i=3$, the actual head transitive on $D$,
and the literal crossing $(6\ b)\delta$ fixing $D\setminus\{6\}$.
It excludes this case without any rank estimate.

It remains that $B$ is imprimitive. Choose any nontrivial
block system, of block size $a\ge2$ and count $z\ge2$.
The accepted two-block theorem excludes $z=2$ in every critical
degree $n\ge18$. The accepted uniform pair-block theorem
excludes $a=2$ in that same range. If $a,z\ge3$, the
accepted remaining-block-size theorem excludes the system:
its forest bound leaves only $(3,4),(4,3)$ and its parity
argument excludes both. These are all possible block systems.
Thus every transitive alternative is impossible.

Finally the acquired exact early-cut calculation says that
colour $3$ has a unique eligible edge if and only if $B$ is
transitive on $E$. Each colour has at least one eligible edge,
since its full omission is intransitive and $G$ is transitive.
The present exclusion therefore gives at least two.
\end{proof}

\paragraph{Attribution and scope.}
This is an assembly of the complete accepted primitive bound,
pure-head obstruction and three newly reviewed imprimitive
proofs. Their source attributions, retained inputs, full common
theories, programs and distinct reviews remain paid. No new
induction hypothesis is used inside the transitive-$B$ exclusion;
the induction that established the critical configuration remains
explicitly inherited. This closes the transitive-$B$ branch only.
Opposite heads for $n\ge26$, the unshared singleton endpoint,
the odd-degree separator refinement and the full L044 recovery
remain outstanding. No whole-original, formal or question credit
is claimed.
\end{document}
